\documentclass[a4paper,twoside]{article}
\usepackage[T1]{fontenc}
\usepackage[bahasa]{babel}
\usepackage{graphicx}
\usepackage{graphics}
\usepackage{float}
\usepackage[cm]{fullpage}
\pagestyle{myheadings}
\usepackage{etoolbox}
\usepackage{setspace} 
\usepackage{lipsum} 
\setlength{\headsep}{30pt}
\usepackage[inner=2cm,outer=2.5cm,top=2.5cm,bottom=2cm]{geometry} %margin
% \pagestyle{empty}

\makeatletter
\renewcommand{\@maketitle} {\begin{center} {\LARGE \textbf{ \textsc{\@title}} \par} \bigskip {\large \textbf{\textsc{\@author}} }\end{center} }
\renewcommand{\thispagestyle}[1]{}
\markright{\textbf{\textsc{Laporan Perkembangan Pengerjaan Skripsi\textemdash Sem. Genap 2015/2016}}}

\onehalfspacing
 
\begin{document}

\title{\@judultopik}
\author{\nama \textendash \@npm} 

%ISILAH DATA BERIKUT INI:
\newcommand{\nama}{Christofer Setiawan}
\newcommand{\@npm}{6182301076}
\newcommand{\tanggal}{24/11/2026} %Tanggal pembuatan dokumen
\newcommand{\@judultopik}{Pengembangan Fitur Ujian dan Download Assignment Pada SharIF-Judge} % Judul/topik anda
\newcommand{\kodetopik}{PAN6102BCS}
\newcommand{\jumpemb}{1} % Jumlah pembimbing, 1 atau 2
\newcommand{\pembA}{Pascal~Alfadian,~Nugroho,~M.Comp.}
\newcommand{\pembB}{-}
\newcommand{\semesterPertama}{61 - Ganjil 26/27} % semester pertama kali topik diambil, angka 1 dimulai dari sem Ganjil 96/97
\newcommand{\lamaSkripsi}{1} % Jumlah semester untuk mengerjakan skripsi s.d. dokumen ini dibuat
\newcommand{\kulPertama}{Skripsi 1} % Kuliah dimana topik ini diambil pertama kali
\newcommand{\tipePR}{B} % tipe progress report :
% A : dokumen pendukung untuk pengambilan ke-2 di Skripsi 1
% B : dokumen untuk reviewer pada presentasi dan review Skripsi 1
% C : dokumen pendukung untuk pengambilan ke-2 di Skripsi 2

% Dokumen hasil template ini harus dicetak bolak-balik !!!!

\maketitle

\pagenumbering{arabic}

\section{Data Skripsi} %TIDAK PERLU MENGUBAH BAGIAN INI !!!
Pembimbing utama/tunggal: {\bf \pembA}\\
Pembimbing pendamping: {\bf \pembB}\\
Kode Topik : {\bf \kodetopik}\\
Topik ini sudah dikerjakan selama : {\bf \lamaSkripsi} semester\\
Pengambilan pertama kali topik ini pada : Semester {\bf \semesterPertama} \\
Pengambilan pertama kali topik ini di kuliah : {\bf \kulPertama} \\
Tipe Laporan : {\bf \tipePR} -
\ifdefstring{\tipePR}{A}{
			Dokumen pendukung untuk {\BF pengambilan ke-2 di Skripsi 1} }
		{
		\ifdefstring{\tipePR}{B} {
				Dokumen untuk reviewer pada presentasi dan {\bf review Skripsi 1}}
			{	Dokumen pendukung untuk {\bf pengambilan ke-2 di Skripsi 2}}
		}
		
\section{Latar Belakang}
	\textit{Coding} merupakan bagian penting bagi seorang mahasiswa informatika yang kelak menjadi seorang programmer. Cara agar seorang mahasiswa informatika bisa memiliki kemampuan programming yang baik adalah dengan latihan dan mengerjakan tugas yang diberikan oleh dosen. Suatu kode dinilai baik atau buruk berdasarkan keakuratan masukan, keluaran dan seberapa efisien kode dihasilkan. Namun setiap orang memiliki cara berpikir dan implementasi kode yang berbeda sehingga penilaian menjadi cukup sulit karena dosen harus menilai kode setiap mahasiswa secara manual.
    
    \textit{Website} SharIF-Judge bekerja dengan memasukkan masukan yang sudah ke dalam kode yang sedang diperiksa. Setelah itu hasil keluaran dari kode tersebut akan dibandingkan dengan keluaran yang sudah disiapkan oleh pembuat soal, penilaian kode didasarkan seberapa akurat keluaran yang dihasilkan. Arsitektur dari \textit{website} SharIF-Judge memakai bahasa pemrograman PHP dan framework codeIgniter lalu untuk basis data memakai \textit{PostgreSQL}. \textit{website} SharIF-Judge  \footnote{https://github.com/ifunpar/SharIF-Judge}ini dipakai oleh jurusan teknik informatika dari Universitas Katolik Parahyangan untuk memudahkan dalam menilai kode yang sudah dikerjakan oleh para mahasiswa. \textit{Website} SharIF-Judge yang digunakan merupakan SharIF-Judge yang sudah dimodifikasi oleh Stillmen Vallian \footnote{Stillmen Vallian, “Kustomisasi Sharif Judge untuk kebutuhan Program Studi Teknik Informatika”, 2018} yang berasal dari ShariF-Judge milik Mohammad Javad Naderi\footnote{Nicholas Aditya Halim, “Implementasi Editor Kode pada SharIF Judge”, 2022}.

	Saat ini SharIF-Judge masih memerlukan dua \textit{instance} terpisah untuk melakukan pengerjaan latihan dan \textit{quiz}, hal ini dilakukan agar ketika mahasiswa sedang ujian atau \textit{quiz} tidak bisa melihat kode atau catatan pada \textit{website} latihan. Namun dua \textit{instance} terpisah ini berdampak pada kurang efisiennya dalam pemeliharaan \textit{website} dan ketika masa \textit{quiz} atau ujian mahasiswa harus mengetik ulang secara manual untuk bisa masuk ke \textit{website} yang dikhususkan untuk \textit{quiz}. Selain itu, tidak ada fitur bagi dosen untuk mengunduh semua \textit{assignment} secara bersamaan sehingga dosen harus mengunduh secara manual.

	Pada tugas akhir ini akan menambahkan dua fitur utama. Fitur pertama adalah pembatasan akses berdasarkan \textit{IPv4} dan penanda \textit{quiz} dalam bentuk \textit{flag}. Fitur ini akan membantu dalam pembatasan lokasi akses jaringan dan penguncian \textit{quiz} dalam satu instance. Lalu akan dikembangkan indikator waktu pengerjaan berbasis visual dan status visual berupa emoji. Fitur kedua, mengunduh seluruh berkas assignment dalam format zip dengan penamaan \texttt{\{id-assignment\}-\{nama-assignment\}-\{nama-soal\}.zip}, hal ini untuk memudahkan bagi para dosen untuk mengecek jawaban para mahasiswa tanpa harus melihat secara manual.
		
\section{Rumusan Masalah}
Rumusan Masalah yang akan dibahas pada tugas akhir adalah sebagai berikut:

	\begin{enumerate}
     \item Bagaimana mengimplementasikan pembatasan akses berdasarkan \textit{IPv4} dan penanda \textit{quiz} dalam satu \textit{instance}?
     \item Bagaimana merancang perubahan tampilan antarmuka agar memiliki status visual berupa emoji dan durasi pengerjaan setiap \textit{assignment} diberikan warna sesuai dengan status \textit{assignment}?
     \item Bagaimana cara melakukan pengunduhan berkas \textit{assignment} yang sudah dikompresi dalam bentuk zip dan memiliki penamaan \texttt{\{id-assignment\}-\{nama-assignment\}-\{nama-soal\}.zip}.
 	\end{enumerate}

\section{Tujuan}

Tujuan yang ingin dicapai tugas akhir ini sebagai berikut:
	\begin{enumerate}
     \item Mengimplementasikan pembatasan akses berdasarkan \textit{IPv4} dan penanda \textit{quiz} dalam satu \textit{instance}.
     \item Mengimplementasikan perubahan tampilan antarmuka agar memiliki status visual berupa emoji dan durasi pengerjaan setiap \textit{assignment} diberikan warna sesuai dengan status \textit{assignment}.
     \item Mengimplementasikan pengunduhan berkas \textit{assignment} yang sudah dikompresi menjadi bentuk zip dan memiliki penamaan \texttt{\{id-assignment\}-\{nama-assignment\}-\{nama-soal\}.zip}.
 	\end{enumerate}

\section{Detail Perkembangan Pengerjaan Skripsi}
Detail bagian pekerjaan skripsi sesuai dengan rencan kerja/laporan perkembangan terkahir :
	\begin{enumerate}
		\item \textbf{Mempelajari PHP}\\
		{\bf Status :} Ada sejak rencana kerja skripsi.\\
		{\bf Hasil :} Survei sudah dilakukan sebanyak 3x pada tanggal X, Y dan Z. Pada kunjungan pertama, diperhaitkan denah museum dan dibuat sketsa berdasarkan pengamatan. Pada kunjunga kedua, bertemu dengan bagian humas museum dan berhasil mendapatkan denah serta melakukan wawancara. Hasil wawancara sudah dibuat dan ada softcopy-nya. Pada kunjungan ketiga, secara khusus dilihat perilaku pengunjung. Masih direncakanan 2x kunjungan lagi. Bukti-bukti kunjungan dapat dilihat di lampiran
		
		\item \textbf{Mempelajari framework CodeIgniter untuk PHP.}\\
		{\bf Status :} Ada sejak rencana kerja skripsi.\\
		{\bf Hasil :}

		\item \textbf{Melakukan studi mengenai cara pembatasan akses oleh IPv4 memakai PHP}\\
		{\bf Status :} Ada sejak rencana kerja skripsi.\\
		{\bf Hasil :}

		\item \textbf{Melakukan studi mengenai mengunduh dan melakukan kompressi file menjadi zip}\\
		{\bf Status :} Ada sejak rencana kerja skripsi.\\
		{\bf Hasil :}

		\item \textbf{Merencanakan perubahan pada struktur dan database SharIF-Judge UNPAR.}\\
		{\bf Status :} Ada sejak rencana kerja skripsi.\\
		{\bf Hasil :}

		\item \textbf{Menulis dokumen skripsi}\\
		{\bf Status :} Ada sejak rencana kerja skripsi.\\
		{\bf Hasil :} 

	\end{enumerate}

\section{Pencapaian Rencana Kerja}
Langkah-langkah kerja yang berhasil diselesaikan dalam Skripsi 1 ini adalah sebagai berikut:
\begin{enumerate}
\item
\item
\item
\end{enumerate}



\section{Kendala yang Dihadapi}
%TULISKAN BAGIAN INI JIKA DOKUMEN ANDA TIPE A ATAU C
Kendala - kendala yang dihadapi selama mengerjakan skripsi :
\begin{itemize}
	\item 
	\item 
	\item 
	\item 
\end{itemize}

\vspace{1cm}
\centering Bandung, \tanggal\\
\vspace{2cm} \nama \\ 
\vspace{1cm}
\begin{center}
    \includegraphics[width=3.5cm, height=2cm, keepaspectratio]{SIgn.png} \\
\end{center}
\vspace{1cm}
\vspace{1cm}

Menyetujui, \\
\ifdefstring{\jumpemb}{2}{
\vspace{1.5cm}
\begin{centering} Menyetujui,\\ \end{centering} \vspace{0.75cm}
\begin{minipage}[b]{0.45\linewidth}
% \centering Bandung, \makebox[0.5cm]{\hrulefill}/\makebox[0.5cm]{\hrulefill}/2013 \\
\vspace{2cm} Nama: \pembA \\ Pembimbing Utama
\end{minipage} \hspace{0.5cm}
\begin{minipage}[b]{0.45\linewidth}
% \centering Bandung, \makebox[0.5cm]{\hrulefill}/\makebox[0.5cm]{\hrulefill}/2013\\
\vspace{2cm} Nama: \pembB \\ Pembimbing Pendamping
\end{minipage}
\vspace{0.5cm}
}{
% \centering Bandung, \makebox[0.5cm]{\hrulefill}/\makebox[0.5cm]{\hrulefill}/2013\\
\vspace{2cm} Nama: \pembA \\ Pembimbing Tunggal
}
\end{document}

